\documentclass{article}

\usepackage[space]{cite}
\usepackage{hyperref}

\title{\bf Citations with the number-only Scheme}
\date{}

\renewcommand{\citepunct}{, and\,}
\renewcommand{\citemid}{;~}
\renewcommand\citeleft{(}
\renewcommand\citeright{)} 
\renewcommand\citedash{\,:\,}

\hypersetup{
	colorlinks=true, citecolor=blue
}

\begin{document}
	\maketitle
	
	\section{Citations Management}
	
	\subsection{Basic Work}
	For more information refer to \cite[P.15]{LGC97}, or 
	refer to \cite[pp\,128\,--130]{GUD97}. \\[3pt]
	You can also have multiple cites \cite{LGC97,GUD97}.
	
	\subsection{Automatic Sorting and Compression}
	\textit{\bfseries Stopped by [nosort, nocompress] 				package options.} \\
	The {\bf cite} package sorts and compresses three 			or more consecutive citations into range format
	\cite{TCM,LGC97,GUD97} % [3,1,2] -> [1,2,3] -> [1-3]
	without sorting or compression it will be [3,1,2].
	
	\subsection{Plain Citing}
	Plain citation with {\bf cite} package is like see
	\citen{TCM}
	
	\pagebreak
	
	\begin{thebibliography}{111}
		%%% The First Ref %%%
		\bibitem{LGC97} 
			Goossens, M., S.~Rahtz, and F.~Mittelbach 					(1997). 
			\newblock \emph{The \LaTeX{} Graphics 						Companion: Illustrating Documents with 
			\TeX{} and PostScript}. 
			\newblock Reading, MA, USA:
			Addison-Wesley Longman.
		%%% The Second Ref %%%
		\bibitem{GUD97} 
			Goossens, M., B.~User, J.~Doe (1997). 
			\newblock Ambiguous citations.
		%%%% Third Ref %%%%
		\bibitem{TCM} 
			D.E.Knuth, Typesetting Concrete 							Mathematics. 
			\newblock \emph{TUGboat}, Apr.\,1989.
	\end{thebibliography}

\end{document}